\section{反向采样是在拿步数换质量}
\label{sec:14-Deep-Learning/08-Diffusion-Models/03}

你把采样步数从 \(50\) 降到 \(10\)，出图快了五倍，整体构图看着还行，但人脸开始出问题：眼睛一高一低，手指数目不对，直线的边缘带着轻微的波纹。第一反应是模型没训好，可同一个权重文件用 \(250\) 步跑，这些毛病一个都没有。权重没变，训练目标没变，变的只有你把那条链切成了多少段。

问题不在模型，在求解器。训练结束时你手上的东西是一个去噪回归器，等价地是一族得分估计——那是一个\emph{向量场}，不是一个采样算法。把向量场变成样本要沿着它走完整条路，而每一步走多远由你决定，走出来的误差会一路累积到终点。这一节要说清楚这笔账怎么算。

\subsection{随机采样：把反向链一步步走完}

最直接的做法是把上一节推出来的反向条件分布拿来直接用。既然 \(p_\theta(x_{t-1}\given x_t)=\mathcal N(\mu_\theta(x_t,t),\sigma_t^2 I)\)，那就从 \(x_T\sim\mathcal N(0,I)\) 出发，逐步采样：

\[
x_{t-1}=\frac{1}{\sqrt{\alpha_t}}\left(x_t-\frac{\beta_t}{\sqrt{1-\bar\alpha_t}}\,\varepsilon_\theta(x_t,t)\right)+\sigma_t z_t,
\qquad z_t\sim\mathcal N(0,I),
\]

最后一步（\(t=1\)）不加噪声。这就是 DDPM 采样。括号里那一项是往条件均值走，\(\sigma_t z_t\) 是每一步重新注入的随机性。

为什么还要往回加噪声？如果每一步都精确地落在条件均值上，你得到的不是分布的一个样本，而是一串条件期望的复合——那是一个把质量往高密度区压缩的映射，跑到终点会塌成模糊的``平均脸''。条件均值只是这一小步的中心，实际的 \(x_{t-1}\) 应该在它周围按 \(\sigma_t\) 散开；这份散开是分布本身的宽度，不是误差。换个角度说，这个迭代是在数值模拟一个反向时间的随机微分方程，\(\sigma_t z_t\) 就是其中的扩散项，\hyperref[sec:06-Random-Processes/06-Diffusion-and-SDE/04]{``数值模拟与时间反演的直觉''}讲过这类模拟的基本结构。

一千步走下来质量很好，问题是慢。每一步都要跑一次完整的网络前向，出一张图就是一千次前向。而这一千步里的绝大部分，看上去只是在做非常小的修正。能不能少走几步？

\subsection{去掉随机项，轨迹就变成了曲线}

DDIM 给出的答案是：先把随机性拿掉。它的出发点是一个不太显然的观察——训练目标 \(L_{\mathrm{simple}}\) 只依赖边缘分布 \(q(x_t\given x_0)\)，不依赖前向过程是不是马尔可夫的。于是可以构造一族别的前向过程，它们的每一个 \(q(x_t\given x_0)\) 都和原来完全一样（所以训好的网络照用不误），但反向过程的随机性可以连续调节，一直调到零。取零的那一端，更新规则是

\[
x_{s}=\sqrt{\bar\alpha_{s}}\;\underbrace{\frac{x_t-\sqrt{1-\bar\alpha_t}\,\varepsilon_\theta(x_t,t)}{\sqrt{\bar\alpha_t}}}_{\hat x_0(x_t,t)}\;+\;\sqrt{1-\bar\alpha_{s}}\,\varepsilon_\theta(x_t,t),
\qquad s<t .
\]

读法很清楚：先用当前的噪声预测反推出一个干净样本的估计 \(\hat x_0\)，再按目标时刻 \(s\) 的调度系数把它重新加噪回去。整个式子里没有新抽的随机数，给定 \(x_T\) 就唯一确定了整条轨迹。

关键在于 \(s\) 和 \(t\) 不必相邻。上式只用到 \(\bar\alpha_t\) 和 \(\bar\alpha_s\) 两个数，任取一个递减的时刻子序列 \(T=t_N>t_{N-1}>\dots>t_0=0\) 都能跑。从一千步里挑五十步出来，算法照常工作。这就是跳步。

\subsection{为什么跳步是合法的：它在解一个 ODE}

确定性轨迹这四个字不只是``省掉了随机数''，它是一个结构性的事实，也是本节最值得记住的连接。

对同一族边缘分布 \(\{p_t\}\)，除了那个随机微分方程之外，还存在一个\emph{确定性}的常微分方程，其解轨迹在每个时刻的分布也恰好是 \(p_t\)。它叫概率流 ODE：

\[
\frac{dx}{dt}=-\tfrac12\beta(t)\Bigl(x+\nabla_x\log p_t(x)\Bigr)
=-\tfrac12\beta(t)\left(x-\frac{\varepsilon_\theta(x,t)}{\sqrt{1-\bar\alpha_t}}\right).
\]

同一组边缘分布，两条不同的路径：SDE 版本每条轨迹都在随机游走，ODE 版本每条轨迹都是一条光滑曲线，但两者在任一时刻的``粒子分布''是同一个。生成样本只要求终点的分布对，不要求路径长什么样，所以两条都能用。

而一旦承认采样是在解 ODE，第 8 部分那套东西就整个搬过来了，一个概念对一个概念：

采样步数就是积分步数，时刻子序列就是步长序列，跳步就是把步长调大。DDIM 的更新规则在做一次变量替换之后正好是这个 ODE 的显式 Euler 步，所以它继承 Euler 的全部性质——\hyperref[sec:08-Numerical-Analysis/06-ODE/01]{``Euler 方法：从局部斜率到全局误差''}算过这笔账：局部截断误差是 \(O(h^2)\)，走 \(1/h\) 步累积成 \(O(h)\) 的全局误差，一阶。步数减半，终点误差大致翻倍。你把 \(50\) 步压到 \(10\) 步，误差涨了五倍，涨出来的部分表现为终点偏离了真正的数据流形，而``偏离数据流形''在图像上就是眼睛不对称、手指数目不对——那些需要长程一致性才成立的结构，正是最先崩掉的东西。

既然是一阶方法的通病，解法也是现成的。\hyperref[sec:08-Numerical-Analysis/06-ODE/02]{``Runge--Kutta 与自适应步长''}的核心结论是：在一步之内多算几次斜率、按合适的权重组合，可以把局部误差压到 \(O(h^{p+1})\)，于是同样的精度只需要少得多的步数。扩散采样里那些被称作``高阶求解器''的东西做的就是这件事，只是它们还多用了一层结构——上面那个 ODE 是\emph{半线性}的，右端 \(-\frac12\beta x\) 是线性项，只有含 \(\varepsilon_\theta\) 的那一项是非线性的。线性部分可以精确积分（就是一个指数因子），不需要数值近似；把它解析地解掉、只对非线性部分做多项式外插，得到的指数积分器在同样的函数调用次数下精度高出一大截。这是数值分析里的老办法用在了新地方，不是扩散模型特有的发明。

刚性也照样出现。\hyperref[sec:08-Numerical-Analysis/06-ODE/03]{``稳定区域与刚性方程''}说过，显式方法的步长上限由右端函数的 Lipschitz 常数决定，右端变化越剧烈，步子越迈不动。这个 ODE 的右端含得分场 \(\nabla\log p_t\)，而 \(t\to0\) 时 \(p_t\) 逼近数据分布，密度在流形附近急剧变化，得分场的 Lipschitz 常数随之爆掉。所以低噪声那一端天然是刚性的，步长必须收紧。这解释了实践中一条看似琐碎的经验：时刻子序列不能均匀取，要在低噪声端加密。这条经验的来源是刚性，不是玄学。

\subsection{条件对账：跳步靠的是哪一条}

跳步的合法性押在``轨迹是某个 ODE 的解''这一条上，别的条件都可以松，这一条不能。

DDPM 每一步注入一份新的独立噪声，走出来的是随机微分方程的一条实现，不是任何常微分方程的解。它没有一条确定的曲线可供你``用更大的步长去积分''——把两步并成一步，你并没有沿同一条曲线走得更远，你是换了一个被模拟的过程。后果不只是精度差一点：Euler--Maruyama 这类格式对 SDE 的收敛阶本来就低于 ODE 情形的对应格式，而更麻烦的是，在大步长下每步注入的噪声量与那一步应有的分布宽度对不上，产生的偏差不是``细化步长就会消失''的那种数值误差，是模拟对象本身错了。所以同样压到十步，DDPM 的退化比 DDIM 严重得多。

反过来说，随机性也不全是负担。SDE 采样每一步注入的噪声有一点纠错作用：如果前面某一步把样本推到了低密度区，接下来的漂移项会把它拉回来，而注入的扰动让它有机会离开一个不该待的位置。ODE 采样没有这个机制，一旦某一步走歪，误差就沿着确定的轨迹一路带到终点。所以在步数充足（几百步）的区间里，SDE 采样的样本质量常常反而略好于 ODE 采样。这是经验观察，不同数据集和调度下强弱关系会变。步数少的时候则毫无悬念是 ODE 采样占优，因为那时主导误差的是离散化而不是轨迹漂移。

\begin{center}
\begin{tikzpicture}[font=\small,>=stealth]
% ===== 坐标轴 =====
\draw[->,black!55] (-0.2,0) -- (9.9,0) node[right,font=\footnotesize] {步数};
\draw[->,black!55] (-0.2,0) -- (-0.2,3.45);
\node[rotate=90,font=\footnotesize] at (-0.85,1.7) {样本质量};
% 对数横轴刻度（手工放置）
\foreach \p/\lab in {0/1, 0.903/2, 2.097/5, 3.0/10, 3.903/20, 5.097/50, 6.0/100, 9.0/1000}{
  \draw[black!45] (\p,0) -- (\p,-0.12);
  \node[below,font=\footnotesize] at (\p,-0.12) {\lab};
}
% 上限线
\draw[black!35,dashed] (-0.2,3.0) -- (9.7,3.0);
\node[right,font=\footnotesize,black!55] at (7.4,3.22) {该权重能达到的上限};
% ===== 三条曲线 =====
\draw[green!45!black,thick,smooth] plot coordinates
  {(0,0.28)(0.903,0.95)(2.097,2.10)(3.0,2.66)(3.903,2.88)(5.097,2.95)(6.0,2.97)(7.19,2.98)(9.0,2.99)};
\draw[blue!70,thick,smooth] plot coordinates
  {(0,0.15)(0.903,0.52)(2.097,1.38)(3.0,2.02)(3.903,2.56)(5.097,2.86)(6.0,2.93)(7.19,2.96)(9.0,2.98)};
\draw[red!70,thick,smooth] plot coordinates
  {(0,0.04)(0.903,0.13)(2.097,0.42)(3.0,0.88)(3.903,1.52)(5.097,2.30)(6.0,2.72)(7.19,2.90)(9.0,2.98)};
% ===== 拐点 =====
\fill[green!45!black] (3.24,2.74) circle (2.1pt);
\draw[green!45!black,densely dotted] (3.24,2.74) -- (3.24,-0.55);
\node[below,font=\footnotesize,green!35!black] at (3.24,-0.55) {\(\approx\)10 步};
\fill[blue!70] (4.98,2.84) circle (2.1pt);
\draw[blue!70,densely dotted] (4.98,2.84) -- (4.98,-0.55);
\node[below,font=\footnotesize,blue!70] at (4.98,-0.55) {\(\approx\)50 步};
\fill[red!70] (6.90,2.88) circle (2.1pt);
\draw[red!70,densely dotted] (6.90,2.88) -- (6.90,-0.55);
\node[below,font=\footnotesize,red!70] at (6.90,-0.55) {\(\approx\)200 步};
% ===== 曲线标注 =====
\node[right,font=\footnotesize,green!35!black] at (0.55,2.62) {高阶求解器};
\node[right,font=\footnotesize,blue!70] at (3.15,1.62) {DDIM（一阶确定性）};
\node[right,font=\footnotesize,red!70] at (4.35,1.02) {DDPM（随机采样）};
\end{tikzpicture}

\emph{三条曲线的高度上限是同一条虚线——求解器换不来更好的模型，它们争的只是多快爬到那条线。要记住的是三个拐点的相对位置：去掉随机性把拐点左移一次，提高积分阶数再左移一次。横轴是对数刻度，每一次左移省下的都是数倍的计算量。}

\end{center}

图里三个拐点的具体位置是示意的。步数和质量之间没有普适的定量关系，只有这条经验曲线的\emph{形状}——先陡后平、最终饱和——在各种设置下都稳定出现。拐点落在哪里取决于数据集、噪声调度、时刻子序列的取法、求解器的阶数，甚至取决于你用哪个指标衡量``质量''：按感知指标看已经饱和的地方，按似然看可能还在爬。没有任何定理告诉你多少步够，这是一条必须自己在自己的设置上测出来的曲线。

\subsection{蒸馏换掉的不是求解器，是模型}

高阶求解器能把拐点左移，但左移不到头。它们逼近的是同一个 ODE 的精确解，而那条解轨迹本身有多长、多弯，是模型定的；把步数压到三步、两步，任何求解器都会失效，因为轨迹的曲率在那个尺度上已经不能用几次斜率求值刻画了。

蒸馏走的是另一条路，它不再试图更准地解那个 ODE，而是训练\emph{另一个模型}，让它一步就完成教师模型走两步（或者更多步）的效果。渐进式蒸馏反复做这件事：教师用 \(2N\) 步，学生学会用 \(N\) 步达到同样的终点，学生再当教师，如此折半，几轮下来步数从上千压到个位数。一致性模型更激进，直接学一个把轨迹上任意一点映到轨迹端点的函数，训练时用的约束是``同一条轨迹上的不同点必须映到同一个输出''。

这个区分要说清楚：求解器是在\emph{使用}训练好的得分场，蒸馏是在\emph{学一个新函数}。所以蒸馏要额外的训练成本，而且学生的输出分布不再等于教师的分布，只是它的一个近似——压到一两步时，样本多样性通常明显下降，因为一个几乎无损的一步映射要把整个高斯先验准确地搬到数据分布上，难度和原本那个``一次困难的分布匹配''已经很接近了。整个单元开头拆解掉的那件难事，在极端压缩下又长了回来。这条路上省下的采样成本，是拿训练成本和多样性换来的。

至此，从纯噪声走回数据的机器已经完整：一个训练好的去噪回归器，加上一个把它的输出积分成轨迹的求解器。但它现在只会生成``像训练数据的东西''，你没法告诉它你想要什么。下一节把控制信号接进来，而接入的位置恰好就是这一节反复出现的那个量——得分。
